% answers_short.tex —— 1988 年·全国卷·理科 · 简短答案速查
% 由 tools/make_short_answers.py 自 answers.tex 的 \gksolhead/\gkanswer 提取，
% 并与 1988/1988_数学真题_全国理_答案要点.md 逐题比对后生成，勿手改（改源 MD 后重跑）。

\begin{gknotice}
\textbf{使用说明}：本节只给最终答案便于快速核对，完整推导见解析版。答案为编写者独立推导，\textbf{非官方标准答案}。
\end{gknotice}

\gkpart{题目}
\begin{itemize}[leftmargin=2.6em]
  \item[\textbf{1.}] B
  \item[\textbf{2.}] C
  \item[\textbf{3.}] B
  \item[\textbf{4.}] A
  \item[\textbf{5.}] D
  \item[\textbf{6.}] A
  \item[\textbf{7.}] C
  \item[\textbf{8.}] D
  \item[\textbf{9.}] C
  \item[\textbf{10.}] D
  \item[\textbf{11.}] B
  \item[\textbf{12.}] B
  \item[\textbf{13.}] D
  \item[\textbf{14.}] B
  \item[\textbf{15.}] A
  \item[\textbf{16.}] 模为 $2$，辐角主值为 $-\dfrac\pi6$
  \item[\textbf{17.}] $x=-2$
  \item[\textbf{18.}] $\tan\dfrac\theta2=-3$
  \item[\textbf{19.}] $\sin\alpha=\dfrac1{\sqrt5}$
  \item[\textbf{20.}] $\dfrac1{q^5}$
  \item[\textbf{21.}] $\dfrac{a(a^2+3)}2$
  \item[\textbf{22.}] $\dfrac{\sqrt3\pi a^3}{36}$
  \item[\textbf{23.}] 若 $a>1$，则 $\dfrac12\log_a t\leq\log_a\dfrac{t+1}{2}$；若 $0<a<1$，不等号反向；$t=1$ 时等号成立。
  \item[\textbf{24.}] (1) 经过图象上任意两个不同点的直线不平行于 $x$ 轴；(2) 函数图象关于直线 $y=x$ 成轴对称图形
  \item[\textbf{25.}] $0<p<\dfrac13$
\end{itemize}
